By \claimref{std-algebra.product-decomposition}, $a \rfloor X = \tfrac12(aX - \hat{X} a)$. Since the grade involution is an automorphism (\claimref{std-algebra.involution-properties}), $\tfrac12\bigl(aXY - \hat{X}\hat{Y}a\bigr) = \tfrac12(aX - \hat{X}a)Y + \tfrac12\hat{X}(aY - \hat{Y}a)$, which is the first identity.

For the outer product, take an orthogonal basis and $a = e_i$, $X = e_K$. If $i \notin K$, $e_i \rfloor e_K = 0$; if $i \in K$, say in position $m$ of $K$, then $e_K = (-1)^{m-1} e_i e_{K \setminus \{i\}}$ and $e_i \rfloor e_K = e_i e_K = (-1)^{m-1} Q(e_i)\, e_{K \setminus \{i\}}$, by \claimref{std-algebra.product-decomposition}. This is the interior product $\iota_{B(e_i, \cdot)}$ applied to the basis element $e_K$ of $\Lambda V$, through $\Phi$ (\claimref{std-algebra.exterior-isomorphism}). The interior product with a linear form is a graded derivation of $\wedge$ that squares to zero and takes the value $B(a, b)$ on a vector $b$; this gives the second and third identities, and the expansion follows from the second by induction on $k$.
